\chapter[2008 --- Prizes]{2008: Multiple Prizes}
\label{ch:2008_Hausen}

\section*{Main Contribution}

\textit{Discovered human papillomaviruses (HPV) as the cause of cervical cancer, leading to vaccines that can prevent this deadly disease.}\cite{nobel-2008,lecture-2008}

\section{Official Citation}

for his discovery of human papillomaviruses causing cervical cancer; and for their discovery of human immunodeficiency virus

\section{Harald zur Hausen}

\subsection{Background and Historical Context}

Harald zur Hausen, a German virologist born in 1936 in Gelsenkirchen, made a discovery that fundamentally changed the understanding of cancer causation and led to the development of one of the major cancer prevention tools in history. His identification of human papillomaviruses (HPV) as the causative agents of cervical cancer challenged long-held assumptions about the relationship between viruses and cancer, and ultimately led to the development of vaccines that have the potential to virtually eliminate one of the most common and deadly cancers affecting women worldwide.

The idea that viruses might cause cancer was not new when zur Hausen began his research in the 1970s. By that time, several animal viruses had been shown to cause tumors in their natural hosts, including the Rous sarcoma virus in chickens and the Epstein-Barr virus in humans, which was linked to Burkitt lymphoma and nasopharyngeal carcinoma. However, the viral etiology of cervical cancer---the second most common cancer in women worldwide---remained unproven. The search for a viral cause of cervical cancer had been ongoing since the 1960s, when the oncologist Johannes Gissmann and the virologist Harald zur Hausen's colleague Lutz Gissmann first suggested that herpes simplex virus type 2 (HSV-2) might be involved.

In the 1970s, zur Hausen and his colleagues at the University of Erlangen-Nuremberg began systematic investigations into the viral etiology of cervical cancer. While initial efforts to detect HSV-2 in cervical tumor samples using molecular hybridization techniques were largely unsuccessful, zur Hausen noticed something intriguing: when DNA from cervical cancer biopsies was hybridized with probes for HSV-2, the signal was inconsistent and often absent. This observation led zur Hausen to hypothesize that a different virus, rather than HSV-2, might be responsible for cervical cancer. He directed his attention toward human papillomaviruses, a diverse group of DNA viruses known to cause warts and other benign skin lesions.

\subsection{The Discovery/Contribution}

Zur Hausen's hypothesis was that specific types of HPV, particularly those associated with genital warts, might be the true causative agents of cervical cancer. In the late 1970s and early 1980s, his laboratory developed novel molecular techniques for detecting HPV DNA in clinical samples. Using a method called recombinant cloning, zur Hausen and his collaborators isolated DNA fragments from cervical cancer biopsies and screened them for the presence of HPV sequences.

In 1983, zur Hausen and his colleague Lutz Gissmann reported the discovery of HPV type 16 (HPV-16) DNA in cervical cancer specimens. This was followed in 1984 by the identification of HPV type 18 (HPV-18). These were the first definitive demonstrations that specific HPV types were present in cervical cancer tissue, providing the strongest evidence to date for a viral etiology of the disease. Subsequent studies revealed that HPV-16 and HPV-18 were detected in approximately 70\% of cervical cancer specimens worldwide, with additional high-risk HPV types (including HPV-31, HPV-33, and HPV-45) accounting for most of the remaining cases.

Zur Hausen's work demonstrated that HPV infection was not merely associated with cervical cancer but was causally linked to the disease. The key evidence for causation came from several lines of investigation. First, epidemiological studies showed that women infected with high-risk HPV types had a dramatically increased risk of developing cervical cancer compared to uninfected women. Second, molecular studies revealed that the HPV genes E6 and E7, which encode oncoproteins that inactivate the tumor suppressors p53 and retinoblastoma protein (pRb), respectively, were expressed in virtually all HPV-positive cervical cancers. Third, experimental studies showed that HPV DNA could transform normal cells into cancerous cells in laboratory culture systems.

The identification of HPV as the causative agent of cervical cancer had immediate and far-reaching implications for cancer prevention. It enabled the development of screening tests that could detect HPV infection in women, allowing early identification of those at risk for cervical cancer. More importantly, it opened the door to the development of prophylactic vaccines that could prevent HPV infection and thereby prevent cervical cancer.

\subsection{Impact on Modern Medicine}

The impact of zur Hausen's discovery on modern medicine has been transformative. The development of prophylactic HPV vaccines, beginning with the approval of Gardasil by the U.S. FDA in 2006 and Cervarix in 2009, has provided a powerful tool for preventing HPV infection and cervical cancer. Clinical trials demonstrated that these vaccines were highly effective in preventing infection with the HPV types they targeted, with efficacy rates exceeding 90\% for the prevention of HPV-16 and HPV-18 infections and associated cervical precancerous lesions.

Widespread HPV vaccination programs have been implemented in many countries, and early data from these programs have shown dramatic reductions in HPV infection rates, genital warts, and cervical precancerous lesions among vaccinated populations. In countries with high vaccination coverage, such as Australia, which was among the first to implement a national HPV vaccination program in 2007, the results have been particularly striking, with near-elimination of HPV-16 and HPV-18 infections in vaccinated cohorts.

Beyond cervical cancer, HPV has been linked to several other types of cancer, including vulvar, vaginal, anal, penile, and oropharyngeal cancers. The HPV vaccines protect against these HPV-related cancers as well, expanding the potential benefits of vaccination. Zur Hausen's discovery has thus had a broad impact on cancer prevention across multiple organ systems.

Harald zur Hausen was awarded the Nobel Prize in Physiology or Medicine in 2008 ``for his discovery of human papillomaviruses causing cervical cancer.'' His work stands as one of the major examples of how basic research in virology can lead to the development of life-saving medical interventions.

\section{Fran\c{c}oise Barr\'e-Sinoussi}

\subsection{Background and Historical Context}

Fran\c{c}oise Barr\'e-Sinoussi, a French virologist born in 1947 in Paris, played a central role in the identification of the human immunodeficiency virus (HIV) as the causative agent of acquired immunodeficiency syndrome (AIDS). Together with her colleague Luc Montagnier, she made a discovery that would transform the understanding of one of the most devastating pandemics in human history and open the door to the development of diagnostic tests, antiretroviral therapies, and ultimately a path toward controlling the HIV/AIDS epidemic.

The AIDS epidemic emerged in the early 1980s as one of the most alarming public health crises in modern history. In 1981, the U.S. Centers for Disease Control and Prevention (CDC) reported an unusual cluster of cases of \textit{Pneumocystis carinii} pneumonia and Kaposi's sarcoma among young gay men in Los Angeles and New York City. These diseases, which were normally seen only in individuals with severely compromised immune systems, suggested that a new and devastating form of immunodeficiency was affecting previously healthy young people. The epidemic spread rapidly, particularly among gay men, intravenous drug users, and recipients of blood transfusions, and the cause of the disease remained unknown.

The search for the causative agent of AIDS was one of the most intense and competitive research efforts in the history of virology. In the United States, Robert Gallo and his colleagues at the National Cancer Institute were pursuing the hypothesis that a retrovirus related to human T-cell leukemia virus (HTLV) might be responsible. In France, a team at the Pasteur Institute in Paris, led by Luc Montagnier with Barr\'e-Sinoussi as a key collaborator, was conducting independent investigations.

\subsection{The Discovery/Contribution}

In January 1983, Montagnier and Barr\'e-Sinoussi received lymph node biopsies from a patient with AIDS-related complex (ARC), a precursor to full-blown AIDS. Barr\'e-Sinoussi, who had expertise in retrovirus culture, attempted to isolate a virus from these samples. Using a technique called co-culture, in which patient cells were cultured with stimulated peripheral blood lymphocytes, she successfully isolated a novel retrovirus from the lymph node tissue.

The virus, initially called lymphadenopathy-associated virus (LAV), exhibited several distinctive features. Unlike HTLV, which infected primarily CD4+ T cells and promoted their proliferation, the new virus appeared to be cytopathic---it killed the cells it infected. Barr\'e-Sinoussi and Montagnier observed that the virus caused the formation of multinucleated giant cells (syncytia) in infected cultures, a hallmark of viral cytopathic effect. They also demonstrated that the virus was tropic for CD4+ T cells, the key helper cells of the immune system.

The identification of the virus as the cause of AIDS required several additional steps. Barr\'e-Sinoussi and Montagnier isolated the virus from multiple AIDS patients and demonstrated that it was consistently present in individuals with AIDS or AIDS-related conditions. They also showed that the virus could be transmitted to healthy cells in culture and that this transmission was associated with the destruction of CD4+ T cells, providing a direct link between viral infection and the immunodeficiency seen in AIDS patients.

The race to identify the cause of AIDS was marked by intense competition and controversy, particularly between the French and American teams. Robert Gallo's group published a paper in \textit{Science} in May 1984 claiming to have discovered the virus, which they called HTLV-III, and to have developed a blood test for it. However, it was subsequently shown that Gallo's isolate was derived from a sample provided by Montagnier's group. After years of dispute, the controversy was resolved in 1987 through an agreement brokered by the U.S. and French governments, which shared credit for the discovery and arranged for the sharing of royalties from blood test patents.

\subsection{Impact on Modern Medicine}

The identification of HIV as the causative agent of AIDS had an immediate and transformative impact on medicine and public health. It enabled the development of diagnostic blood tests that could screen blood supplies for HIV infection, dramatically reducing the risk of transfusion-transmitted HIV. It also enabled epidemiological tracking of the epidemic, revealing the global extent of HIV infection and identifying populations at notable risk.

More broadly, the identification of HIV opened the door to the development of antiretroviral therapies. The first anti-HIV drug, zidovudine (AZT), was approved in 1987, followed by a succession of increasingly effective medications. The development of highly active antiretroviral therapy (HAART) in the mid-1990s, which combined multiple drugs targeting different stages of the viral life cycle, transformed HIV/AIDS from a death sentence into a manageable chronic condition for those with access to treatment.

Fran\c{c}oise Barr\'e-Sinoussi was awarded the Nobel Prize in Physiology or Medicine in 2008, shared with Luc Montagnier, ``for their discovery of human immunodeficiency virus.'' Her work has been recognized as a landmark achievement in virology and public health, and her continued advocacy for HIV/AIDS research and treatment has maintained the visibility of the epidemic and the importance of ongoing efforts to develop a vaccine and cure.

\section{Luc Montagnier}

\subsection{Background and Historical Context}

Luc Montagnier, a French virologist born in 1932 in Chabris, France, shared the Nobel Prize with Fran\c{c}oise Barr\'e-Sinoussi for his role in the discovery of HIV. Montagnier led the team at the Pasteur Institute that first isolated the virus and characterized its basic properties, work that was essential for understanding the cause of AIDS and developing the tools to combat the epidemic.

Montagnier's interest in retroviruses dated back to the 1970s, when he studied interferon and the effects of viral infections on the immune system. When the AIDS epidemic emerged in the early 1980s, Montagnier recognized that the disease might be caused by a retrovirus and directed his laboratory's efforts toward identifying such a virus. His experience with retrovirus isolation techniques and his leadership of a talented team of researchers, including Barr\'e-Sinoussi, positioned his laboratory to make the critical discovery.

\subsection{The Discovery/Contribution}

Montagnier's leadership and scientific vision were essential to the success of the HIV discovery effort. He made the crucial decision to focus on lymph node biopsies from patients with AIDS-related conditions, rather than blood samples, which were being studied by other groups. This decision proved to be pivotal, as the lymph nodes of patients with early-stage disease were more likely to contain replicating virus than peripheral blood.

Montagnier also played an important role in the characterization of the virus after its initial isolation by Barr\'e-Sinoussi. He contributed to the demonstration that HIV was a lentivirus---a type of retrovirus associated with chronic, slowly progressive diseases---and to the understanding of its genetic diversity. Montagnier's group was among the first to sequence the HIV genome and to identify the major structural and regulatory genes of the virus.

In the years following the discovery of HIV, Montagnier continued to make important contributions to HIV/AIDS research. He studied the mechanisms of viral pathogenesis, the immune response to HIV infection, and the potential for vaccine development. He also became a prominent advocate for increased investment in HIV/AIDS research and for equitable access to antiretroviral therapy in developing countries.

In later years, Montagnier became a controversial figure due to his public statements questioning the safety of vaccines and promoting alternative theories about the origins of HIV. These positions were not supported by the scientific community and diverged significantly from the mainstream scientific understanding of HIV/AIDS. Despite these controversies, Montagnier's contribution to the discovery of HIV remains a landmark achievement in the history of medicine.

\subsection{Impact on Modern Medicine}

The discovery of HIV by Montagnier and Barr\'e-Sinoussi has had a lasting impact on global health. The development of antiretroviral therapy has saved millions of lives and transformed the HIV/AIDS pandemic from a rapidly fatal disease to a manageable chronic condition in many parts of the world. However, the epidemic continues to pose a major public health challenge, with approximately 38 million people living with HIV worldwide and approximately 1.5 million new infections occurring each year.

Montagnier's work also contributed to the broader understanding of retroviral pathogenesis and the mechanisms by which viruses can cause chronic disease and immune dysfunction. The study of HIV has informed the understanding of other viral infections and has contributed to the development of antiviral therapies and vaccine strategies for a range of infectious diseases.

Luc Montagnier shared the Nobel Prize in Physiology or Medicine in 2008 with Fran\c{c}oise Barr\'e-Sinoussi ``for their discovery of human immunodeficiency virus.'' Their work represents one of the major achievements in the history of medicine, providing the foundation for the global response to one of the most devastating epidemics in human history.


\section{Modern Developments}

zur Hausen, Barré-Sinoussi, and Montagnier's virus work has led to modern cancer prevention and HIV treatment. HPV vaccines (Gardasil, Cervarix) have dramatically reduced cervical cancer rates. HIV antiretroviral therapy has transformed HIV from a death sentence to a chronic condition. Understanding of viral oncogenesis has informed development of oncolytic virus therapy. Understanding of HIV reservoirs has led to research on functional cures for HIV.


\section{Bibliography}

\begin{thebibliography}{9}
\bibitem{nobel-2008}
Nobel Prize Outreach, ``The Nobel Prize in Physiology or Medicine 2008,''\emph{NobelPrize.org}.\\
\url{https://www.nobelprize.org/prizes/medicine/2008/summary/}

\bibitem{lecture-2008}
Harald zur Hausen, ``Nobel Lecture,''\emph{NobelPrize.org}. Winner-authored primary source; downloadable from the official Nobel Prize archive.\\
\url{https://www.nobelprize.org/prizes/medicine/2008/hausen/lecture/}
\end{thebibliography}
